%! TEX root = ../note-dqm-transform.tex

\section{Introduction}
\label{sec--intro}

Let \((\mathbf{X}, \mathcal{B})\) be a measurable space.
We use linear functional notation for integrals: given a measure \(m\) on
\((\mathbf{X}, \mathcal{B})\), \(m [f] := \int f (x) \; m (\mathrm{d} x)\).
Furthermore, denote the Euclidean norm by \(|\cdot|\), i.e.
\(|x|^{2} = \sum_{j = 1}^{d} x_{j}^{2}\) for any \(x \in \mathbb{R}^{d}\).

Let \(\mu\) be a \(\sigma\)-finite measure on \((\mathbf{X}, \mathcal{B})\) and
let \(\Theta\) be a non-empty subset of \(\mathbb{R}^{d}\) for fixed \(d \in
\mathbb{N}\).
Let \(\mathbf{Q} = \left\{ Q_{\tau} : \tau \in \Theta \right\}\) be a family
of probability measures dominated by \(\mu\).
Denote the \(\mu\)-densities and their square-roots by
\begin{equation}
  q_{\tau} (\cdot) := q (\cdot; \tau) := \frac{\mathrm{d}
  Q_{\tau}}{\mathrm{d} \mu}, \quad \text{and} \quad g_{\tau} (\cdot) := g
  (\cdot; \tau) := \sqrt{q (\cdot; \tau)}.
  \label{eqn--q-density-and-sqrt}
\end{equation}
Since \(q_{\tau}\) is a \(\mu\)-density for a probability measure, \(\int
g_{\tau}^{2} \; \mathrm{d} \mu = \int q_{\tau} \; \mathrm{d} \mu = 1\).
Hence \(\tau \mapsto g_{\tau}\) maps \(\Theta\) into \(\mathscr{L}_{2} (\mu)\).

\begin{definition}[Hellinger differentiability]
\label{def--hellinger-diff}
The map \(\tau \mapsto Q_{\tau}\) is Hellinger differentiable at \(\theta \in
\Theta\) if there is a \(d\)-vector of \(\mathscr{L}_{2} (\mu)\) functions
\(\dot{g}_{\theta} (\cdot) = \dot{g} \left( \cdot; \theta \right)\) such that
\begin{equation}
  \mu \left[ \rho_{\tau, \theta}^{2} \right] = o \left( \left| \tau -
  \theta \right|^{2} \right) \text{ as } \tau \to \theta \text{ while }
  \tau \in \Theta, \quad \text{where} \quad
  \rho_{\tau, \theta} = g_{\tau} - g_{\theta} - \left(
  \tau - \theta \right)^{\prime} \dot{g}_{\theta}.
  \label{eqn--hellinger-diff-def-remainder}
\end{equation}
The map \(\tau \mapsto Q_{\tau}\) is Hellinger differentiable if it is
Hellinger differentiable at every \(\tau \in \Theta\).
\end{definition}

Let \(T : (\mathbf{X}, \mathcal{B}) \to \left( \mathbf{X}^{\prime},
\mathcal{B}^{\prime} \right)\) be a measurable transformation and
consider the collection of pushforward probability measures:
\begin{equation*}
  \mathbf{P} = \{ P_{\tau} = Q_{\tau} \circ T^{- 1} : \tau \in \Theta\}.
\end{equation*}
The question of preservation of Hellinger differentiability under \(T\)
essentially boils down to two questions:
is the family \(\mathbf{P}\) dominated by a measure, say \(\nu\), and are the
associated root-densities Hellinger differentiable (against \(\mathscr{L}_{2}
(\nu)\)).
A natural candidate for a dominating measure for \(\mathbf{P}\) would be the
pushforward \(\nu = \mu \circ T^{- 1}\), but this is not necessarily
\(\sigma\)-finite as illustrated by \Cref{eg--non-sigma-finite-pushforward}.

\begin{example}
\label{eg--non-sigma-finite-pushforward}
Let \(T : \mathbb{R}^{2} \to \mathbb{R}^{2}\) be defined by \(T (x, y) = (y,
\mathbf{1} \{x \leq y\})\) and let \(\mu\) be Lebesgue measure on
\(\mathbb{R}^{2}\).
Then \(\nu = \mu \circ T^{- 1}\) is not \(\sigma\)-finite.
A deeper look at this failure of \(\sigma\)-finiteness is given in
\Cref{sec--prf--eg--non-sigma-finite-pushforward}.
\end{example}

If \(\mu\) is a probability measure, its pushforward under any measurable
transformation is automatically \(\sigma\)-finite.
\Cref{thm--sigma-finite-meas-abs-bicont-prob-meas} shows that for
\(\sigma\)-finite \(\mu\), it is always possible to find a probability
measure \(\lambda\) such that both \(\mu \ll \lambda\) and \(\lambda \ll \mu\).
Furthermore, \Cref{thm--hellinger-diff-unaffected-by-dom} shows that
replacing \(\mu\) by \(\lambda\) such that \(\mu \ll \lambda\) does
not affect Hellinger differentiability.


\begin{theorem}
\label{thm--sigma-finite-meas-abs-bicont-prob-meas}
Let \(\mu\) be a positive \(\sigma\)-finite measure on \((\mathbf{X},
\mathcal{B})\).
There is a probability measure \(\lambda\) on \((\mathbf{X}, \mathcal{B})\) such
that \(\mu \ll \lambda\) and \(\lambda \ll \mu\).
\end{theorem}

\begin{proof}
See \Cref{sec--prf--thm--sigma-finite-meas-abs-bicont-prob-meas}.
\end{proof}

\begin{theorem}
\label{thm--hellinger-diff-unaffected-by-dom}
Let \(\mu, \lambda\) be \(\sigma\)-finite measures on \((\mathbf{X},
\mathcal{B})\) such that \(\mu \ll \lambda\).
Denote \(f^{2} = \mathrm{d} \mu / \mathrm{d} \lambda\).
Let \(\mathbf{Q} = \left\{ Q_{\tau} : \tau \in \Theta \right\}\) be a family
of probability measures dominated by \(\mu\) and denote \(g_{\tau}^{2} =
\mathrm{d} Q_{\tau} / \mathrm{d} \mu\).
Suppose also that \(\tau \mapsto g_{\tau}\) is Hellinger differentiable at
\(\theta\) with derivative \(\dot{g}_{\theta}\).
Then \(\mathbf{Q}\) satisfies the following:
\begin{enumerate}[label=(\alph*)]
  \item \label{thm--hellinger-diff-unaffected-by-dom-rn-chain-rule}
    \(\mathbf{Q}\) is dominated by \(\lambda\) and its \(\lambda\)-densities
    are \(g_{\ast, \tau}^{2} := g_{\tau}^{2} f^{2}\);
  \item \label{thm--hellinger-diff-unaffected-by-dom-derivative-def}
    The map \(\tau \mapsto g_{\ast, \tau}\) is Hellinger differentiable at
    \(\theta\) with derivative \(\dot{g}_{\ast, \theta} =
    \dot{g}_{\theta} f\).
\end{enumerate}
\end{theorem}

\begin{proof}
See \Cref{sec--prf--thm--hellinger-diff-unaffected-by-dom}.
\end{proof}

Given
\Cref{thm--sigma-finite-meas-abs-bicont-prob-meas,%
thm--hellinger-diff-unaffected-by-dom}, we can assume that \(\mu\) is a
probability measure.
Then the pushforward measure \(\nu = \mu \circ T^{- 1}\) is also a
probability measure and hence \(\sigma\)-finite.
Let \(\mu [f | T]\) denote the conditional expectation of \(f\) given \(T\).
\Cref{thm--density-preservation} shows that \(P_{\tau} = Q_{\tau} \circ T^{-
1}\) is dominated by \(\nu\) and gives a form for the density \(p_{\tau}\) as a
conditional expectation.

\begin{theorem}
\label{thm--density-preservation}
Suppose that \(Q\) and \(\mu\) are probability measures on \((\mathbf{X},
\mathcal{B})\) such that \(Q \ll \mu\).
Let \(T : \mathbf{X} \mapsto \mathbf{X}^{\prime}\) be measurable from
\((\mathbf{X}, \mathcal{B})\) to \(\left( \mathbf{X}^{\prime},
\mathcal{B}^{\prime} \right)\), and denote \(\nu := \mu \circ T^{- 1}\).
Then \(P = Q \circ T^{- 1} \ll \nu\) and
\begin{equation}
  \frac{\mathrm{d} P}{\mathrm{d} \nu} (t) = \mu \left[ \frac{\mathrm{d}
  Q}{\mathrm{d} \mu} \middle| T = t \right] \quad \nu\text{-a.e}.
  \label{eqn--density-preservation}
\end{equation}
\end{theorem}

\begin{proof}
See \Cref{sec--prf--thm--density-preservation}.
\end{proof}

%%% Local Variables:
%%% mode: LaTeX
%%% TeX-master: "../note-dqm-transform"
%%% End:
